% Extracted from p05-curriculum
\subsection{Curriculum Strategies for TinyStories}
\label{ch-curriculum}

\paragraph{Background}

\paragraph{Curriculum learning}

\citet{bengio2009curriculum} introduced curriculum learning by analogy
with human teaching: present easy examples first, then gradually harder
ones. Follow-up work formalised this as a pacing function $c(t)$, the
fraction of the (difficulty-sorted) data available at training step $t$,
growing from a small $c(0)$ to $c(T) = 1$. Reported benefits cluster in
two regimes: faster early convergence, and better final performance on
tasks with a clear difficulty structure. Reported failures cluster where
difficulty is ill-defined or where the pacing interacts badly with the
learning-rate schedule. Both the promise and the failure modes motivate
preregistration: the difficulty metric and pacing must be fixed before
training, or the experimenter can always find \emph{some} ordering that
``worked'' after the fact.

\paragraph{TinyStories as a testbed}

TinyStories contains short children's stories generated by larger models
under tight vocabulary and narrative constraints \citep{eldan2023tinystories}.
Three properties matter for our purposes. First, stories vary genuinely
in difficulty: some use only the simplest words and one character, others
use rarer vocabulary and multi-step plots. Second, the corpus is fully
available offline, so corpus-wide statistics (unigram frequencies, length
distributions) are computable exactly. Third, the reference 44M run gives
us a calibrated baseline: we know what shuffled training achieves at this
budget, so curriculum conditions have a fixed target to beat.

\paragraph{Story-complexity metrics}

We define five metrics, each computable from the story text alone plus
corpus-wide unigram counts. All are defined before any training.

\paragraph{Vocabulary rarity.} For a story $S$ with tokens $w_1, \dots,
w_n$ and corpus unigram probabilities $p(\cdot)$,
{\small
\begin{equation}
R(S) = -\frac{1}{n}\sum_{i=1}^{n} \log p(w_i).
\end{equation}}
Stories built from rare words score high. This is the mean unigram
surprisal, cheap and stable.

\paragraph{Mean sentence length.} Longer sentences in children's stories
correlate with subordinate clauses and heavier syntax. Computed in
tokens after the same BPE tokenisation used in training.

\paragraph{Narrative depth.} A proxy for plot complexity:
{\small
\begin{equation}
\begin{aligned}
N(S) = \#\text{distinct named characters} \\
+\; \#\text{causal/temporal connectives},
\end{aligned}
\end{equation}}
where connectives are drawn from a fixed list (\emph{because, so, then,
after, before, when, but}). Crude, but transparent and reproducible.

\paragraph{Lexical diversity.} Type--token ratio over the story. High
diversity at fixed length indicates richer word choice.

\paragraph{Story length.} Total tokens. Longer stories demand longer-range
coherence.

\paragraph{Difficulty score.} Each metric is z-scored across the corpus,
and the difficulty score is the unweighted mean of the five z-scores:
\begin{equation}
D(S) = \frac{1}{5}\sum_{m=1}^{5} z_m(S).
\end{equation}
Equal weights are a deliberate, preregistered choice. Tuning weights
after seeing results would convert the study into a fishing expedition;
the weights are therefore fixed before any training run begins.

\paragraph{Proposed experimental protocol}

\paragraph{Fixed factors}

Model, tokeniser, optimiser, learning-rate schedule, batch size, and
total token budget (120M) are identical to the reference run in every
condition. Each condition is trained from three random seeds. Only the
order in which stories are presented differs.

\paragraph{Conditions}

\begin{enumerate}
\item[(a)] \textbf{Shuffled baseline.} Uniform random order, as in the
reference run.
\item[(b)] \textbf{Easy-to-hard, linear pacing.} Stories sorted by
$D(S)$; pacing $c(t) = c_0 + (1 - c_0)\, t/T$ with $c_0 = 0.1$.
\item[(c)] \textbf{Easy-to-hard, root pacing.} Same ordering, but
$c(t) = \sqrt{(1-c_0^2)\, t/T + c_0^2}$. This schedule is concave:
it rises steeply at the start, so harder material is introduced
\emph{earlier} than under linear pacing---testing whether early
exposure to difficulty accelerates learning rather than destabilising it.
\item[(d)] \textbf{Anti-curriculum.} Hard-to-easy with linear pacing.
A control: if \emph{any} non-random ordering helped, this condition
would catch it.
\item[(e)] \textbf{Two-stage.} First half of training on the easiest
50\% of stories (shuffled within), second half on the full corpus
shuffled. Tests whether a brief ``foundation'' phase suffices.
\end{enumerate}

\paragraph{Evaluation}

We track training loss continuously and evaluate held-out perplexity at
eight evenly spaced checkpoints plus the final model, on a fixed
validation split never used for pacing decisions. Beyond perplexity, we
use two targeted probes. A \emph{grammaticality probe}: synthetic
minimal pairs (correct vs.\ agreement-violated sentences in the
TinyStories vocabulary) scored by relative likelihood. A
\emph{narrative-coherence probe}: fixed story beginnings completed by
each model and rated for coherence, either by human raters or by a
fixed LLM judge with a published rubric. Perplexity measures fit;
the probes measure what the fit bought.

\paragraph{Hypotheses}

\begin{itemize}
\item[\textbf{H1.}] Easy-to-hard conditions (b, c) reach lower final
held-out perplexity than the shuffled baseline (a) at the fixed token
budget.
\item[\textbf{H2.}] Root pacing (c) outperforms linear pacing (b): earlier
exposure to harder material accelerates representation learning rather
than destabilising early training.
\item[\textbf{H3.}] The anti-curriculum (d) underperforms the baseline
(a), confirming that ordering effects are directional rather than mere
artefacts of non-random presentation.
\item[\textbf{H4.}] Curriculum gains concentrate in the narrative-coherence
probe rather than raw perplexity: ordering teaches structure more than
it teaches token statistics.
\end{itemize}

Each hypothesis is falsifiable as stated. In particular, if (b) and (c)
fail to beat (a), the honest conclusion is that curricula are not worth
the engineering at this scale---a useful negative result.

\paragraph{Discussion}

\paragraph{Why this testbed can settle the question.} Most curriculum
studies inherit an ambiguity: is the difficulty metric real, or did the
experimenter choose one that flatters the method? Here the metrics are
mechanical, the weights preregistered, and the corpus synthetic enough
that the metrics track something genuine. A null result here would be
evidence, not noise.

\paragraph{Confounds.} Two deserve attention. First, pacing interacts
with the learning-rate schedule: easy data early plus a warmup schedule
may double-count the same effect. We keep the schedule fixed across
conditions precisely so this interaction, if present, is held constant
rather than tuned away. Second, difficulty sorting changes the token
distribution seen early in training, which changes early gradient
statistics; three seeds per condition give a rough sense of the variance
this introduces.

\paragraph{What success would mean.} A clean H1--H2 win would argue for
curriculum as a free improvement to small-model pretraining: no extra
tokens, no extra parameters, only ordering. It would also motivate
testing whether the effect persists at larger budgets, where the
relative cost of sorting is negligible.

\paragraph{Limitations}

This paper proposes; it does not execute. Until the runs exist, the
hypotheses are untested. TinyStories is synthetic, so any finding may
not transfer to natural corpora where difficulty is harder to define.
The difficulty score's equal weights are arbitrary by construction---a
principled weighting would itself require a study. Finally, all
conditions share one architecture; curriculum effects may interact with
model structure (notably the recurrent loops of our reference model) in
ways a single-architecture study cannot see. The three-seed design
proposed here has limited power to separate nearby curricula: with only
three runs per condition, small but real ordering effects may not reach
significance, so a null result should be read as inconclusive rather
than as evidence of no effect.

\paragraph{Conclusion}

We have specified a complete, preregistered experimental programme for
testing curriculum learning in TinyStories pretraining at the 44M
scale: five computable difficulty metrics, one fixed difficulty score,
five training conditions, and four falsifiable hypotheses. The protocol
is designed so that both positive and null results are informative. The
next step is execution.
